\documentclass[12pt, a4paper]{article}

% ── Idioma y codificación ──────────────────────────────────────────────────
\usepackage[utf8]{inputenc}
\usepackage[T1]{fontenc}
\usepackage[spanish, es-nodecimaldot]{babel}

% ── Matemáticas ───────────────────────────────────────────────────────────
\usepackage{amsmath, amssymb, amsthm}
\usepackage{mathtools}
\usepackage{physics}

% ── Fuente ────────────────────────────────────────────────────────────────
\usepackage{lmodern}
\usepackage{microtype}

% ── Geometría de página ───────────────────────────────────────────────────
\usepackage[
  top=2.5cm, bottom=2.5cm,
  left=2.8cm, right=2.8cm
]{geometry}

% ── Encabezado y pie ──────────────────────────────────────────────────────
\usepackage{fancyhdr}
\setlength{\headheight}{15pt}
\pagestyle{fancy}
\fancyhf{}
\fancyhead[L]{\small\nouppercase{\leftmark}}
\fancyhead[R]{\small\thepage}
\renewcommand{\headrulewidth}{0.4pt}

\fancypagestyle{plain}{%
  \fancyhf{}%
  \renewcommand{\headrulewidth}{0pt}%
}

% ── Figuras y diagramas ───────────────────────────────────────────────────
\usepackage{tikz}
\usetikzlibrary{arrows.meta, calc, angles, quotes,
                decorations.pathreplacing,
                decorations.pathmorphing, babel}
\usepackage{pgfplots}
\pgfplotsset{compat=1.18}
\usepackage{float}
\usepackage{caption}
\usepackage{array}
\usepackage{tabularx}
\usepackage{booktabs}
\usepackage{setspace}
\usepackage{enumitem}

% ── Estilo visual global ──────────────────────────────────────────────────
\definecolor{brandblue}{HTML}{000000}
\definecolor{softblue}{HTML}{FFFFFF}
\definecolor{linegray}{HTML}{000000}
\colorlet{blue}{black}
\colorlet{red}{black}
\colorlet{green}{black}
\colorlet{orange}{black}

\captionsetup{
  font=small,
  labelfont={bf,color=brandblue},
  textfont=it,
  justification=centering
}

\newcolumntype{Y}{>{\raggedright\arraybackslash}X}
\newcolumntype{C}{>{\centering\arraybackslash}X}
\newcolumntype{L}[1]{>{\raggedright\arraybackslash}p{#1}}
\newcolumntype{B}[1]{>{\bfseries\raggedright\arraybackslash}p{#1}}
\setlength{\tabcolsep}{7pt}
\renewcommand{\arraystretch}{1.18}

\tikzset{
  every picture/.style={line cap=round, line join=round, >=Stealth},
  every path/.append style={draw=black},
  every node/.append style={
    font=\footnotesize,
    text=black,
    fill=white,
    fill opacity=1,
    text opacity=1,
    inner sep=1.2pt
  }
}

% ── Hipervínculos ─────────────────────────────────────────────────────────
\usepackage[hidelinks]{hyperref}

% ── Cajas para resultados destacados ──────────────────────────────────────
\usepackage{mdframed}
\usepackage[skins,breakable]{tcolorbox}
\mdfdefinestyle{resultado}{%
  linecolor=black,
  backgroundcolor=white,
  linewidth=0.8pt,
  innertopmargin=6pt,
  innerbottommargin=6pt,
  innerleftmargin=8pt,
  innerrightmargin=8pt,
  skipabove=8pt,
  skipbelow=8pt,
}

\tcbset{
  demostracionbox/.style={
    enhanced,
    breakable,
    frame hidden,
    interior hidden,
    borderline west={0.7pt}{0pt}{black},
    left=8pt, right=6pt,
    top=4pt, bottom=4pt,
    before skip=8pt, after skip=8pt,
  }
}

% ── Entornos matemáticos ──────────────────────────────────────────────────
\theoremstyle{plain}
\newtheorem{teorema}{Teorema}[section]
\newtheorem{proposicion}[teorema]{Proposición}
\newtheorem{lema}[teorema]{Lema}
\newtheorem{corolario}[teorema]{Corolario}

\theoremstyle{definition}
\newtheorem{definicion}[teorema]{Definición}
\newtheorem{ejemplo}{Ejemplo}[section]

\theoremstyle{remark}
\newtheorem*{nota}{Nota}
\newtheorem*{observacion}{Observación}

% ── Formato de secciones ──────────────────────────────────────────────────
\usepackage{titlesec}
\usepackage{etoolbox}
\titleformat{\section}
  {\large\bfseries\color{brandblue}}
  {\thesection.}{0.6em}{}
  [\vspace{2pt}\color{linegray}\titlerule\vspace{4pt}]

\titleformat{\subsection}
  {\normalsize\bfseries\color{brandblue!90!black}}
  {\thesubsection.}{0.5em}{}

\titleformat{\subsubsection}
  {\normalsize\bfseries\color{brandblue!75!black}}
  {\thesubsubsection.}{0.5em}{}

\titlespacing*{\section}{0pt}{0pt}{0.75\baselineskip}
\titlespacing*{\subsection}{0pt}{0.95\baselineskip}{0.45\baselineskip}
\titlespacing*{\subsubsection}{0pt}{0.75\baselineskip}{0.35\baselineskip}

\pretocmd{\section}{\clearpage}{}{}

% ── Macros ────────────────────────────────────────────────────────────────
\newcommand{\R}{\mathbb{R}}
\newcommand{\C}{\mathbb{C}}
\newcommand{\N}{\mathbb{N}}
\newcommand{\Z}{\mathbb{Z}}
\newcommand{\eps}{\varepsilon}
\newcommand{\ve}[1]{\vec{#1}}
\newcommand{\uvec}[1]{\hat{#1}}
\DeclareMathOperator{\sen}{sen}

% ── Ritmo de lectura ──────────────────────────────────────────────────────
\raggedbottom
\setstretch{1.08}
\setlength{\parskip}{4pt}
\setlength{\parindent}{0pt}
\setlength{\abovedisplayskip}{7pt}
\setlength{\belowdisplayskip}{7pt}
\setlength{\abovedisplayshortskip}{6pt}
\setlength{\belowdisplayshortskip}{6pt}

\setlist[itemize]{leftmargin=1.35em, itemsep=2pt, topsep=4pt, parsep=0pt, partopsep=0pt}
\setlist[enumerate]{leftmargin=1.55em, itemsep=2pt, topsep=4pt, parsep=0pt, partopsep=0pt}
\setlist[description]{style=nextline, leftmargin=0pt, labelwidth=\linewidth,
  labelindent=0pt, labelsep=0pt, itemindent=0pt, listparindent=0pt,
  itemsep=4pt, topsep=4pt, parsep=0pt, partopsep=0pt}

\renewcommand{\proofname}{Demostración}
\renewcommand{\qedsymbol}{$\square$}
\BeforeBeginEnvironment{proof}{\begin{tcolorbox}[demostracionbox]}
\AfterEndEnvironment{proof}{\end{tcolorbox}}
\apptocmd{\proof}{\small\itshape}{}{}

\newcommand{\formulaname}[1]{\textbf{#1}\par\vspace{2pt}}

% ══════════════════════════════════════════════════════════════════════════
\begin{document}
% ══════════════════════════════════════════════════════════════════════════

% ── PORTADA ───────────────────────────────────────────────────────────────
\begin{titlepage}
  \thispagestyle{plain}
  \centering
  \vspace*{1.6cm}

  \begin{tikzpicture}[remember picture,overlay]
    \draw[line width=0.8pt]
      ($(current page.north west)+(1.6cm,-1.6cm)$)
      rectangle
      ($(current page.south east)+(-1.6cm,1.6cm)$);
  \end{tikzpicture}

  {\Large\scshape Mecánica y Ondas\par}
  \vspace{0.25cm}
  {\large Curso 2025--2026\par}

  \vspace{0.9cm}
  {\Huge\bfseries Apuntes Unificados\par}
  \vspace{0.4cm}
  {\large Teoría, demostraciones y formulario\par}

  \vspace{0.8cm}
  \rule{0.78\linewidth}{0.8pt}
  \vspace{0.6cm}
  {\normalsize
  Tema 1: Oscilaciones libres, amortiguadas y forzadas\\
  Tema 2: Osciladores acoplados, modos normales y cuerda vibrante\par}
  \vspace{0.6cm}
  \rule{0.78\linewidth}{0.8pt}

  \vspace{1.2cm}
  \begin{tikzpicture}[>=Stealth, scale=1.15]
    % Sistema masa-resorte
    \fill[black!15] (-0.25,-0.4) rectangle (0,0.4);
    \draw[thick] (0,-0.4) -- (0,0.4);
    \draw[decorate,
          decoration={coil, aspect=0.5, segment length=4.5pt, amplitude=4pt}]
      (0,0) -- (2.2,0);
    \draw[thick] (2.2,-0.28) rectangle (2.82,0.28);
    \node at (2.51,0) {$m$};
    \node[above=4pt] at (1.1,0.18) {$k$};
    % Flecha x
    \draw[->] (2.82,0) -- (3.4,0) node[right] {$x$};
    % Onda sinusoidal
    \draw[->] (4.0,-0.55) -- (4.0,0.85) node[above] {$x(t)$};
    \draw[->] (4.0,-0.55) -- (7.8,-0.55) node[right] {$t$};
    \draw[domain=0:3.5, samples=120, thick]
      plot (\x+4.0, {0.55*sin(3.14159*\x r) - 0.0});
    \node[right] at (7.1, 0.6) {\small $A\,\sen(\omega_0 t - \delta)$};
  \end{tikzpicture}

  \vfill
  \begin{tabular}{@{}rl@{}}
    \textbf{Autor:}      & Luis López \\
    \textbf{Titulación:} & Grado en Física \\
    \textbf{Fecha:}      & Abril 2026
  \end{tabular}
\end{titlepage}

% ── ÍNDICE ────────────────────────────────────────────────────────────────
\tableofcontents

% ══════════════════════════════════════════════════════════════════════════
\section{Tema 1: Oscilaciones Lineales}
% ══════════════════════════════════════════════════════════════════════════

\subsection{Introducción}

Para una partícula existe una posición de equilibrio estable que tomamos como origen.
Cuando la partícula sufre un desplazamiento desde el origen aparece una fuerza
recuperadora que tiende a devolverla a su posición original. En general esta fuerza es una
función complicada del desplazamiento y quizá de la velocidad; aquí solo consideramos el
caso en que $F = F(x)$.

Expandiendo en serie de Taylor alrededor del punto de equilibrio:
\begin{equation}
  F(x) = F_0 + x\!\left(\frac{dF}{dx}\right)_{\!0}
         + \frac{1}{2}x^2\!\left(\frac{d^2F}{dx^2}\right)_{\!0} + \cdots
\end{equation}

En el equilibrio $F_0 = 0$. Para desplazamientos suficientemente pequeños los términos de
orden superior son despreciables y el primero que sobrevive es el lineal, lo que da la
\textbf{ley de Hooke}:
\begin{mdframed}[style=resultado]
\[
  F(x) = -kx, \qquad k = -\!\left(\frac{dF}{dx}\right)_{\!0} > 0
\]
\end{mdframed}

La condición $k > 0$ garantiza la estabilidad del equilibrio (la fuerza se opone al
desplazamiento).

\subsection{Oscilador armónico simple}

Aplicando la segunda ley de Newton con la fuerza de Hooke:
\[
  F = ma = -kx \implies m\ddot{x} = -kx
\]

Se obtiene la \textbf{ecuación del oscilador armónico simple} (OAS):
\begin{mdframed}[style=resultado]
\[
  \ddot{x} + \omega_0^2\,x = 0, \qquad \omega_0^2 = \frac{k}{m}
\]
\end{mdframed}

donde $\omega_0$ es la \emph{frecuencia propia} o \emph{pulsación natural}. La solución
general es:
\begin{mdframed}[style=resultado]
\[
  x(t) = A\,\sen(\omega_0 t - \delta)
\]
\end{mdframed}

Las constantes $A$ (amplitud) y $\delta$ (fase inicial) se determinan con las condiciones
iniciales $x(0)$ y $\dot{x}(0)$. El movimiento es periódico con período
$T = 2\pi/\omega_0$ y frecuencia $\nu_0 = \omega_0/(2\pi)$.

El movimiento del OAS recibe el nombre de \emph{oscilación libre}: una vez iniciado, no
se detiene nunca, pues la energía mecánica se conserva:
\[
  E = \tfrac{1}{2}kA^2 = \tfrac{1}{2}m\omega_0^2 A^2 = \text{cte.}
\]
Esto es una idealización; en la realidad, las fuerzas disipativas acaban extinguiendo el
movimiento.

\subsection{Oscilaciones amortiguadas}

Introduciendo una fuerza de rozamiento viscoso $-b\dot{x}$ proporcional a la velocidad:
\[
  m\ddot{x} + b\dot{x} + kx = 0
\]

Dividiendo por $m$ y definiendo $\beta = b/(2m)$:
\begin{mdframed}[style=resultado]
\[
  \ddot{x} + 2\beta\dot{x} + \omega_0^2\,x = 0
\]
\end{mdframed}

Proponiendo $x(t)=e^{rt}$ se obtiene la ecuación característica
$r^2 + 2\beta r + \omega_0^2 = 0$, con raíces $r = -\beta \pm \sqrt{\beta^2 - \omega_0^2}$.
La solución general es:
\[
  x(t) = e^{-\beta t}\!\left[A\exp\!\left(\sqrt{\beta^2-\omega_0^2}\,t\right)
         + B\exp\!\left(-\sqrt{\beta^2-\omega_0^2}\,t\right)\right]
\]

Pueden ocurrir tres casos según la relación entre $\omega_0$ y $\beta$:

\subsubsection{Subamortiguado \texorpdfstring{($\omega_0^2 > \beta^2$)}{}}

La raíz del discriminante es imaginaria: $\sqrt{\beta^2-\omega_0^2}=i\Omega$ con
$\Omega=\sqrt{\omega_0^2-\beta^2}$. La solución oscila con amplitud decreciente:
\[
  x(t) = A\,e^{-\beta t}\cos(\Omega t + \varphi)
\]
La frecuencia amortiguada $\Omega < \omega_0$ y la amplitud decae como $e^{-\beta t}$.
El \emph{factor de calidad} $Q = \omega_0/(2\beta)$ mide cuántos ciclos tarda en
amortiguarse.

\subsubsection{Amortiguamiento crítico \texorpdfstring{($\omega_0^2 = \beta^2$)}{}}

Raíz doble $r = -\beta$. La solución general es:
\[
  x(t) = (A + Bt)\,e^{-\beta t}
\]
El sistema retorna al equilibrio lo más rápido posible sin oscilar. Es el régimen óptimo
para suprimir vibraciones.

\subsubsection{Sobreamortiguado \texorpdfstring{($\omega_0^2 < \beta^2$)}{}}

Con $\gamma = \sqrt{\beta^2 - \omega_0^2}$, las dos raíces son reales y negativas:
\[
  x(t) = e^{-\beta t}\!\left[A\,e^{\gamma t} + B\,e^{-\gamma t}\right]
\]
El sistema vuelve al equilibrio sin oscilar, pero más lentamente que en el caso crítico.

\subsection{Oscilaciones forzadas}

El caso más sencillo de oscilación forzada es aquel en que la fuerza externa varía
periódicamente con el tiempo:
\[
  m\ddot{x} + b\dot{x} + kx = F_0\cos\omega t
\]
\begin{mdframed}[style=resultado]
\[
  \ddot{x} + 2\beta\dot{x} + \omega_0^2\,x = A\cos\omega t, \qquad A = \frac{F_0}{m}
\]
\end{mdframed}

La solución general es la suma de la solución homogénea (transitorio, se amortigua) y
una solución particular (régimen permanente):
\[
  x_g(t) = e^{-\beta t}\!\left[A\exp\!\left(\sqrt{\beta^2-\omega_0^2}\,t\right)
            + B\exp\!\left(-\sqrt{\beta^2-\omega_0^2}\,t\right)\right]
\]
\[
  x_p(t) = D\cos(\omega t - \delta)
\]

Sustituyendo $x_p$ en la ecuación diferencial se obtienen la amplitud estacionaria y el
desfase:
\begin{mdframed}[style=resultado]
\[
  D = \frac{F_0/m}{\sqrt{(\omega_0^2-\omega^2)^2 + 4\beta^2\omega^2}},
  \qquad
  \tan\delta = \frac{2\beta\omega}{\omega_0^2 - \omega^2}
\]
\end{mdframed}

Cuando $\omega \to \omega_0$ (y $\beta$ es pequeño), la amplitud $D$ se hace grande:
esto es la \textbf{resonancia}. En ausencia de amortiguamiento ($\beta = 0$) la amplitud
diverge en $\omega = \omega_0$.

\subsection{Ejemplos}

\begin{ejemplo}[Oscilador amortiguado]
Una masa de $1\,\mathrm{kg}$ está unida a un resorte de constante elástica
$k = 25\,\mathrm{N/m}$. El sistema se mueve sobre una superficie horizontal con
coeficiente de amortiguamiento $b = 6\,\mathrm{kg/s}$. La masa se suelta desde el
reposo a $0{,}10\,\mathrm{m}$ de la posición de equilibrio. Determinar la expresión del
desplazamiento.

\medskip
\textit{Parámetros:}
\[
  \omega_0^2 = k/m = 25\,\mathrm{rad^2/s^2},\qquad
  \beta = b/(2m) = 3\,\mathrm{rad/s}
\]
Como $\omega_0^2 = 25 > \beta^2 = 9$, el sistema es subamortiguado:
\[
  \Omega = \sqrt{\omega_0^2 - \beta^2} = 4\,\mathrm{rad/s}
\]
La solución general es $x(t) = e^{-3t}(A\cos 4t + B\sin 4t)$. Aplicando $x(0) = 0{,}10$
y $\dot{x}(0) = 0$:
\[
  A = 0{,}10\,\mathrm{m}, \quad B = \frac{\beta A}{\Omega} = 0{,}075\,\mathrm{m}
\]
\[
  \boxed{x(t) = e^{-3t}\!\left(0{,}10\cos 4t + 0{,}075\sin 4t\right)\,\mathrm{m}}
\]
\end{ejemplo}

\begin{ejemplo}[Oscilador forzado]
Con los mismos datos pero $b = 4\,\mathrm{kg/s}$ y con fuerza periódica
$F(t) = 10\cos(3t)\,\mathrm{N}$ aplicada desde $t=0$. La masa parte del reposo a
$0{,}10\,\mathrm{m}$ del equilibrio. Determinar la expresión del desplazamiento.

\medskip
\textit{Parámetros:}
\[
  \beta = 2\,\mathrm{rad/s},\quad
  \Omega = \sqrt{21}\,\mathrm{rad/s},\quad
  \omega = 3\,\mathrm{rad/s},\quad
  F_0/m = 10\,\mathrm{m/s^2}
\]
Amplitud estacionaria:
\[
  D = \frac{10}{\sqrt{(25-9)^2 + 4\cdot 4\cdot 9}} =
      \frac{10}{\sqrt{256+144}} = \frac{10}{20} = 0{,}5\,\mathrm{m}
\]
Desfase: $\tan\delta = 2\cdot 2\cdot 3/(25-9) = 3/4 \Rightarrow \delta \approx 0{,}6435\,\mathrm{rad}$.

La solución completa es $x(t) = x_g(t) + 0{,}5\cos(3t-\delta)$, donde los coeficientes
del transitorio $x_g$ se determinan aplicando las condiciones iniciales.
\end{ejemplo}

% ══════════════════════════════════════════════════════════════════════════
\section{Tema 2: Osciladores Acoplados}
% ══════════════════════════════════════════════════════════════════════════

\subsection{Introducción}

Los osciladores acoplados son sistemas formados por dos o más osciladores que
interactúan entre sí mediante algún tipo de acoplamiento: fuerzas mecánicas, campos
eléctricos, señales químicas o interacciones biológicas. A diferencia de los osciladores
independientes, el comportamiento de cada oscilador acoplado no solo depende de sus
propiedades individuales, sino también de la influencia mutua que ejercen los demás.

Un sistema físico real es normalmente capaz de vibrar en muchos modos diferentes y puede
resonar a muchas frecuencias. Los tres tipos fundamentales de oscilador que sirven de
base al estudio son:
\begin{itemize}
  \item \textbf{Oscilador armónico simple:}
        $x(t) = A\,\sen(\omega t + \delta)$
  \item \textbf{Oscilador amortiguado:}
        $x(t) = x_0\,e^{-\gamma t}\sen(\Omega t + \delta)$
  \item \textbf{Oscilador forzado:}
        $x(t) = x_0\,e^{-\gamma t}\sen(\Omega t + \delta) + X\cos(\omega t + \delta^*)$
\end{itemize}

\subsection{Dos osciladores acoplados: ecuaciones de movimiento}

Consideremos dos péndulos idénticos de longitud $l$ y masa $m$, unidos por un resorte
de constante $k$ en sus posiciones de equilibrio. Sean $x_A$ y $x_B$ los desplazamientos
horizontales respecto al equilibrio.

Las fuerzas sobre cada masa son la fuerza restauradora del péndulo $(-m\omega_0^2 x)$ y
la del resorte de acoplamiento $(\pm k(x_A-x_B))$, con $\omega_0^2 = g/l$. Aplicando la
segunda ley de Newton:
\begin{align*}
  m\frac{d^2x_A}{dt^2} + m\omega_0^2 x_A + k(x_A - x_B) &= 0 \\[4pt]
  m\frac{d^2x_B}{dt^2} + m\omega_0^2 x_B - k(x_A - x_B) &= 0
\end{align*}

Definiendo $\omega_c^2 = k/m$ (frecuencia de acoplamiento):
\begin{mdframed}[style=resultado]
\begin{align*}
  \frac{d^2x_A}{dt^2} + (\omega_0^2+\omega_c^2)x_A - \omega_c^2 x_B &= 0 \\[4pt]
  \frac{d^2x_B}{dt^2} + (\omega_0^2+\omega_c^2)x_B - \omega_c^2 x_A &= 0
\end{align*}
\end{mdframed}

Estas dos ecuaciones \emph{no pueden resolverse independientemente} por estar acopladas.
En el caso general, A y B oscilan con igual amplitud pero con dos frecuencias distintas
a la vez, lo que produce una modulación visible.

\subsection{Modos normales: consideraciones de simetría}

Cuando el sistema es simétrico existen soluciones especiales en las que ambas masas
oscilan a la misma frecuencia con relación de amplitudes fija: los
\textbf{modos normales de vibración}.

\subsubsection{Modo simétrico (en fase): \texorpdfstring{$\omega'' = \omega_0$}{}}

Si $x_A = x_B$ el resorte de acoplamiento no se deforma y su contribución es nula. Cada
péndulo oscila como si estuviera solo:
\[
  x_A = C\cos\omega_0 t,\qquad x_B = C\cos\omega_0 t,\qquad \omega_0 = \sqrt{g/l}
\]
A y B oscilan con igual amplitud e igual frecuencia.

\subsubsection{Modo antisimétrico (en antifase): \texorpdfstring{$\omega'=\sqrt{\omega_0^2+2\omega_c^2}$}{}}

Si $x_A = -x_B$ el resorte se estira/comprime al máximo y actúa como fuerza restauradora
adicional:
\[
  x_A = C\cos\omega' t,\qquad x_B = -C\cos\omega' t,\qquad
  \omega' = \sqrt{\omega_0^2 + 2\omega_c^2}
\]
A y B oscilan con igual amplitud y desfasadas $180°$. El acoplamiento eleva la frecuencia
por encima de $\omega_0$, por lo que siempre $\omega' > \omega_0$.

\subsection{Coordenadas normales}

El \textbf{método de coordenadas normales} transforma el sistema acoplado en dos
osciladores independientes mediante el cambio de variable:
\begin{mdframed}[style=resultado]
\[
  q_1 = x_A + x_B,\qquad q_2 = x_A - x_B
\]
\end{mdframed}

Sumando y restando las ecuaciones de movimiento:
\begin{align}
  \ddot{q}_1 + \omega_0^2\,q_1 &= 0 \label{eq:q1}\\[4pt]
  \ddot{q}_2 + \omega'^2\,q_2  &= 0,\qquad \omega'^2 = \omega_0^2 + 2\omega_c^2
  \label{eq:q2}
\end{align}

Las coordenadas normales $q_1$, $q_2$ oscilan de forma desacoplada, representando los
\emph{modos normales} del sistema. Las soluciones son:
\[
  q_1 = C\cos\omega_0 t,\qquad q_2 = D\cos\omega' t
\]

Recuperando las coordenadas originales:
\begin{mdframed}[style=resultado]
\[
  x_A = \tfrac{1}{2}(q_1+q_2) =
        \tfrac{1}{2}C\cos\omega_0 t + \tfrac{1}{2}D\cos\omega' t
\]
\[
  x_B = \tfrac{1}{2}(q_1-q_2) =
        \tfrac{1}{2}C\cos\omega_0 t - \tfrac{1}{2}D\cos\omega' t
\]
\end{mdframed}

\begin{itemize}
  \item Si $C = 0$: ambas masas oscilan exclusivamente con $\omega'$ (modo antisimétrico).
  \item Si $D = 0$: ambas masas oscilan exclusivamente con $\omega_0$ (modo simétrico).
\end{itemize}

Cualquier movimiento de los péndulos partiendo del reposo puede describirse como una
combinación lineal de estos dos modos normales.

\subsubsection{Condiciones iniciales: batidos}

Un caso físicamente llamativo se produce cuando $x_A(0)=A_0$, $x_B(0)=0$,
$\dot{x}_A(0)=\dot{x}_B(0)=0$. Aplicando estas condiciones:
\[
  A_0 = \tfrac{1}{2}C + \tfrac{1}{2}D,\quad
  0   = \tfrac{1}{2}C - \tfrac{1}{2}D
  \implies C = D = A_0
\]

Las soluciones se reescriben usando suma-a-producto:
\[
  x_A(t) = A_0\cos\!\left(\frac{\omega'-\omega_0}{2}t\right)
              \cos\!\left(\frac{\omega'+\omega_0}{2}t\right)
\]
\[
  x_B(t) = A_0\sen\!\left(\frac{\omega'-\omega_0}{2}t\right)
              \sen\!\left(\frac{\omega'+\omega_0}{2}t\right)
\]

Si el acoplamiento es débil ($\omega_c \ll \omega_0$), la diferencia
$\omega'-\omega_0 \ll \omega_0$ y se producen \textbf{batidos}: la energía oscila
lentamente de A a B y viceversa con frecuencia de modulación $(\omega'-\omega_0)/2$.

\subsection{Método analítico general}

Para resolver el sistema sin recurrir a la simetría se propone que $x_A$ y $x_B$ pueden
oscilar a la misma frecuencia $\omega$:
\[
  x_A = C\cos\omega t,\qquad x_B = C'\cos\omega t
\]

Sustituyendo en las ecuaciones de movimiento:
\begin{align}
  \text{A:}&\quad (-\omega^2+\omega_0^2+\omega_c^2)C - \omega_c^2 C' = 0
  \label{eq:anaA}\\[4pt]
  \text{B:}&\quad -\omega_c^2 C + (-\omega^2+\omega_0^2+\omega_c^2)C' = 0
  \label{eq:anaB}
\end{align}

De las ecuaciones anteriores se pueden extraer dos expresiones para $C/C'$:
\[
  \frac{C}{C'} = \frac{\omega_c^2}{-\omega^2+\omega_0^2+\omega_c^2}
  \qquad\text{y}\qquad
  \frac{C}{C'} = \frac{-\omega^2+\omega_0^2+\omega_c^2}{\omega_c^2}
\]

Para que el sistema tenga solución no trivial estas dos expresiones deben ser iguales:
\[
  \frac{\omega_c^2}{-\omega^2+\omega_0^2+\omega_c^2}
  = \frac{-\omega^2+\omega_0^2+\omega_c^2}{\omega_c^2}
  \implies
  \left(-\omega^2+\omega_0^2+\omega_c^2\right)^2 = \omega_c^4
\]

Esto da exactamente dos frecuencias normales:
\begin{mdframed}[style=resultado]
\[
  \omega''^2 = \omega_0^2
  \quad\Bigl(\tfrac{C}{C'}=1,\ \text{modo simétrico}\Bigr)
  \qquad
  \omega'^2 = \omega_0^2 + 2\omega_c^2
  \quad\Bigl(\tfrac{C}{C'}=-1,\ \text{modo antisimétrico}\Bigr)
\]
\end{mdframed}

La solución general es:
\[
  x_A = C\cos\omega_0 t + D\cos\omega' t,\qquad
  x_B = C\cos\omega_0 t - D\cos\omega' t
\]

\subsection{Oscilaciones forzadas en sistemas acoplados}

Supongamos que sobre la masa A actúa una fuerza $F(t) = F_0\cos\omega t$. Las ecuaciones
de movimiento son:
\begin{align*}
  m\frac{d^2x_A}{dt^2} + m\omega_0^2 x_A + k(x_A-x_B) &= F_0\cos\omega t \\[4pt]
  m\frac{d^2x_B}{dt^2} + m\omega_0^2 x_B - k(x_A-x_B) &= 0
\end{align*}

En coordenadas normales, con $q_1 = x_A+x_B$ y $q_2 = x_A-x_B$:
\begin{align*}
  \frac{d^2q_1}{dt^2} + \omega_0^2 q_1 &= \frac{F_0}{m}\cos\omega t \\[4pt]
  \frac{d^2q_2}{dt^2} + \omega_0'^2 q_2 &= \frac{F_0}{m}\cos\omega t,
  \qquad \omega_0'^2 = \omega_0^2 + 2\omega_c^2
\end{align*}

Cuando el sistema se impulsa con una frecuencia arbitraria $\omega$ se producen grandes
amplitudes para frecuencias próximas a las frecuencias naturales del sistema, y respuesta
pequeña para frecuencias alejadas de ellas. Las soluciones estacionarias son:
\[
  q_1 = C\cos\omega t,\quad C = \frac{F_0/m}{\omega_0^2-\omega^2};
  \qquad
  q_2 = D\cos\omega t,\quad D = \frac{F_0/m}{\omega_0'^2-\omega^2}
\]

Las amplitudes de desplazamiento resultan:
\begin{mdframed}[style=resultado]
\[
  x_A = A\cos\omega t,\quad A = \tfrac{1}{2}(C+D);
  \qquad
  x_B = B\cos\omega t,\quad B = \tfrac{1}{2}(C-D)
\]
\end{mdframed}

Se producen \textbf{resonancias} en $\omega = \omega_0$ (diverge $C$) y en
$\omega = \omega_0'$ (diverge $D$). El sistema tiene exactamente dos frecuencias de
resonancia, correspondientes a sus dos modos normales.

\begin{ejemplo}[Osciladores acoplados: tres resortes, dos masas]
Considérese el sistema formado por tres resortes idénticos de constante $k$ y dos masas
iguales de valor $2m$, dispuesto como: pared$-k-2m-k-2m-k-$pared.

\textit{(a) Ecuaciones de movimiento:}
\[
  2m\ddot{x}_A + 2kx_A - kx_B = 0,\qquad
  2m\ddot{x}_B + 2kx_B - kx_A = 0
\]

\textit{(b) Frecuencias de oscilación} (con $\omega_0^2 = k/(2m)$, $\omega_c^2 = k/(2m)$):
\[
  \omega'' = \sqrt{\frac{k}{2m}}\quad (\text{modo simétrico}),\qquad
  \omega'  = \sqrt{\frac{3k}{2m}}\quad (\text{modo antisimétrico})
\]

\textit{(c) Si las masas se reducen a $m$ y se quieren mantener las mismas frecuencias},
debe cumplirse $k'/(m) = k/(2m)$, por lo que $k' = k/2$.
\end{ejemplo}

\begin{ejemplo}[Dos péndulos con muelle colgante]
Una partícula de masa $m$ se sujeta a un soporte rígido mediante un muelle de constante
$\kappa$ que cuelga verticalmente en el equilibrio. A esta combinación masa-muelle se
sujeta un oscilador igual, con el muelle de este segundo conectado a la masa del primero.
Calcular las pulsaciones características de las oscilaciones verticales monodimensionales.

Las ecuaciones de movimiento son formalmente iguales a las del sistema de péndulos
acoplados con $\omega_0^2 = \kappa/m$ y $\omega_c^2 = \kappa/m$ (el muelle central actúa
tanto de restaurador del segundo como de acoplamiento):
\[
  \omega'' = \omega_0 = \sqrt{\kappa/m},\qquad
  \omega'  = \sqrt{3\kappa/m}
\]

Cuando una de las masas se mantiene fija, la otra oscila con
$\omega_{\text{fija}} = \sqrt{2\kappa/m}$, que es intermedia entre las dos pulsaciones
normales, como era de esperar.
\end{ejemplo}

\subsection{\texorpdfstring{$N$}{N} osciladores acoplados: la cuerda vibrante discreta}

\subsubsection{Principio de superposición de Bernoulli}

Los Bernoulli (padre e hijo) demostraron que un sistema de $N$ masas tiene exactamente
$N$ modos de vibración. En 1753, Daniel Bernoulli enunció el
\textbf{principio de superposición}: \emph{el movimiento general de un sistema vibrante
puede escribirse como superposición de los modos normales}.

\subsubsection{Ecuación de movimiento}

Consideremos una cuerda elástica flexible sometida a tensión $T$ con $N$ partículas de
masa $m$ equiespaciadas una distancia $l$ entre sí (partículas $p = 1,\ldots,N$, extremos
fijos en $p=0$ y $p=N+1$).

Para desplazamientos transversales pequeños $y_p$, en la aproximación de ángulos pequeños
$(\sen\alpha\approx\tan\alpha)$, la componente longitudinal de la tensión da una
contribución $\frac{1}{2}T(\alpha_1^2-\alpha_2^2)$ que se desprecia, y la fuerza
transversal sobre la partícula $p$ es:
\[
  F_p = -\frac{T}{l}(y_p - y_{p-1}) + \frac{T}{l}(y_{p+1}-y_p)
      = \frac{T}{l}(y_{p+1} - 2y_p + y_{p-1})
\]

La ecuación de movimiento resulta:
\begin{mdframed}[style=resultado]
\[
  \frac{d^2y_p}{dt^2} + 2\omega_0^2\,y_p - \omega_0^2(y_{p+1}+y_{p-1}) = 0,
  \qquad \omega_0^2 = \frac{T}{ml}
\]
\end{mdframed}
con condiciones de contorno $y_0 = y_{N+1} = 0$.

\subsubsection{Cálculo de los \texorpdfstring{$N$}{N} modos normales}

Se proponen soluciones sinusoidales $y_p = A_p\cos\omega t$, $p = 1,2,\ldots,N$.
Sustituyendo en la ecuación de movimiento:
\[
  (-\omega^2+2\omega_0^2)A_p - \omega_0^2(A_{p-1}+A_{p+1}) = 0
  \implies
  \frac{A_{p-1}+A_{p+1}}{A_p} = \frac{-\omega^2+2\omega_0^2}{\omega_0^2}
\]

Se propone $A_p = C\sen(p\theta)$, que satisface
$\frac{A_{p-1}+A_{p+1}}{A_p} = 2\cos\theta$.
La condición de contorno en $p=0$ se cumple automáticamente; la de $p=N+1$ exige:
\[
  A_{N+1} = C\sen\!\bigl((N+1)\theta\bigr) = 0
  \implies (N+1)\theta = n\pi,\quad n = 1,2,3,\ldots
\]

Por tanto:
\begin{mdframed}[style=resultado]
\[
  \theta = \frac{n\pi}{N+1}
  \implies
  A_p = C_n\,\sen\!\left(\frac{np\pi}{N+1}\right)
\]
\end{mdframed}

\subsubsection{Propiedades de los \texorpdfstring{$N$}{N} modos normales}

\begin{itemize}
  \item Los $n$ distintos definen $N$ \textbf{frecuencias normales}:
  \begin{mdframed}[style=resultado]
  \[
    \omega_n = 2\omega_0\,\sen\!\left(\frac{n\pi}{2(N+1)}\right),
    \qquad n = 1, 2, \ldots, N
  \]
  \end{mdframed}

  \item El movimiento de la partícula $p$ en el modo $n$ es:
  \[
    y_{p,n}(t) = C_n\,\sen\!\left(\frac{np\pi}{N+1}\right)\cos\omega_n t
  \]
  El desplazamiento depende tanto de la posición de la partícula ($p$) como del modo
  considerado ($n$).

  \item Existen exactamente $N$ modos normales independientes.
\end{itemize}

Para $n=1$ (modo inferior) todas las partículas se mueven en el mismo sentido con perfil
de medio seno. Para $n=2$ existe un nodo central fijo y las dos mitades oscilan en
antifase. En general, el modo $n$ tiene $n-1$ nodos internos. Los modos con $n > N$
repiten los patrones de los modos menores con signo opuesto.

\subsubsection{Oscilaciones longitudinales}

El mismo formalismo se aplica a desplazamientos longitudinales $\xi_p$:
\[
  \frac{d^2\xi_p}{dt^2} + 2\omega_0^2\,\xi_p - \omega_0^2(\xi_{p+1}+\xi_{p-1}) = 0
\]
Las soluciones y frecuencias son formalmente idénticas a las del caso transversal:
\[
  \xi_p(t) = C_n\,\sen\!\left(\frac{np\pi}{N+1}\right)\cos\omega_n t,
  \qquad
  \omega_n = 2\omega_0\,\sen\!\left(\frac{n\pi}{2(N+1)}\right)
\]

\subsection{Límite continuo: la cuerda vibrante continua}

\subsubsection{Del sistema discreto al continuo}

Una cuerda con ambos extremos fijos tiene un número discreto de estados de vibración
natural bien definidos. Cada punto de la cuerda vibra transversalmente con movimiento
armónico simple de amplitud constante, y la frecuencia es la misma para todas las partes
de la cuerda. Dichas vibraciones representan los modos normales de la cuerda.

Cuando $N\to\infty$ manteniendo fijos la longitud total $L = (N+1)l$ y la masa total
$M = Nm$ (densidad lineal $\mu = M/L$), el sistema discreto converge a una cuerda
continua. Para los modos de baja frecuencia ($n \ll N$) se cumple
$n\pi/[2(N+1)] \ll 1$ y $\sen\theta \approx \theta$, por lo que:
\[
  \omega_n = 2\omega_0\,\sen\!\left(\frac{n\pi}{2(N+1)}\right)
  \xrightarrow{N\to\infty}
  n\frac{\pi}{L}\sqrt{\frac{T}{\mu}} = n\omega_1
\]

Los desplazamientos pasan de forma discreta a forma continua:
\begin{mdframed}[style=resultado]
\[
  y_n(x,t) = C_n\,\sen\!\left(\frac{n\pi x}{L}\right)\cos\omega_n t
\]
\end{mdframed}

La frecuencia máxima del sistema discreto es
$\omega_{\max} = 2\omega_0\sen\!\bigl(N\pi/[2(N+1)]\bigr) \approx 2\omega_0$ para $N$
grande.

\subsubsection{Derivación de la ecuación de ondas}

Supongamos una cuerda de longitud $L$, densidad lineal uniforme $\mu$ y tensión $T$, con
ambos extremos fijos. Para oscilaciones transversales pequeñas $y(x,t)$, considerando un
elemento de longitud $\Delta x$:
\[
  F_y = T\sen(\theta+\Delta\theta) - T\sen\theta \approx T\,\Delta\theta
  \qquad (F_x \approx 0\ \text{para ángulos pequeños})
\]

Dado que $\tan\theta = \partial y/\partial x$, para ángulos pequeños
$\sec^2\theta\,\Delta\theta = (\partial^2 y/\partial x^2)\,\Delta x$, y por tanto
$T\,\Delta\theta \approx T(\partial^2 y/\partial x^2)\,\Delta x$. Aplicando la segunda
ley de Newton al elemento $(\mu\,\Delta x)$:
\[
  T\frac{\partial^2 y}{\partial x^2}\,\Delta x = \mu\,\Delta x\,\frac{\partial^2 y}{\partial t^2}
\]

Se obtiene la \textbf{ecuación de ondas}:
\begin{mdframed}[style=resultado]
\[
  \frac{\partial^2 y}{\partial x^2} = \frac{1}{v^2}\frac{\partial^2 y}{\partial t^2},
  \qquad
  v = \left(\frac{T}{\mu}\right)^{\!1/2}
\]
\end{mdframed}

\subsubsection{Modos normales de la cuerda continua}

Suponiendo soluciones de separación de variables $y(x,t)=f(x)\cos\omega t$:
\[
  \frac{\partial^2 y}{\partial t^2} = -\omega^2 f(x)\cos\omega t,\qquad
  \frac{\partial^2 y}{\partial x^2} = \frac{d^2f}{dx^2}\cos\omega t
  \implies \frac{d^2f}{dx^2} = -\frac{\omega^2}{v^2}f
\]

La solución es $f(x) = A\sen(\omega x/v)$. Las condiciones de contorno $y(0,t)=y(L,t)=0$
imponen:
\[
  \frac{\omega L}{v} = n\pi \implies \omega_n = n\frac{\pi v}{L} = n\omega_1
\]

Las frecuencias y longitudes de onda de los modos normales son:
\begin{mdframed}[style=resultado]
\[
  \nu_n = \frac{nv}{2L} = \frac{n}{2L}\left(\frac{T}{\mu}\right)^{\!1/2},
  \qquad
  \lambda_n = \frac{2L}{n}
\]
\end{mdframed}

El movimiento de la cuerda en el modo $n$ es:
\begin{mdframed}[style=resultado]
\[
  y_n(x,t) = A_n\,\sen\!\left(\frac{n\pi x}{L}\right)\cos\omega_n t,
  \qquad
  \omega_n = \frac{n\pi}{L}\left(\frac{T}{\mu}\right)^{\!1/2} = n\omega_1
\]
\end{mdframed}

Los puntos donde $y_n=0$ para todo $t$ son los \textbf{nodos}
($x = kL/n$, $k=0,1,\ldots,n$); los puntos de máxima amplitud son los
\textbf{antinodos}.

\subsubsection{Cuerda forzada en un extremo}

Cuando un extremo está sometido a una fuerza periódica, las condiciones de contorno
cambian:
\[
  y(0,t) = B\cos\omega t,\qquad y(L,t) = 0
\]

La función espacial adopta la forma $f(x) = A\sen(\omega x/v + \alpha)$. Aplicando las
condiciones de contorno:
\begin{align*}
  \text{En }x=L:&\quad \frac{\omega L}{v} + \alpha = p\pi \\[4pt]
  \text{En }x=0:&\quad B = A\sen\alpha \implies
    A = \frac{B}{\sen\!\left(p\pi - \dfrac{\omega L}{v}\right)}
\end{align*}

La amplitud diverge cuando $\omega L/v = p\pi$, es decir, cuando $\omega = \omega_p$:
la cuerda entra en \textbf{resonancia} con el modo $p$.

\begin{ejemplo}[Cuerda con \texorpdfstring{$N=2$}{N=2} masas]
Se sujeta por sus extremos a dos soportes fijos una cuerda de longitud $3l$ y masa
despreciable, con tensión $T$.

\textit{(a) Una sola masa $m$ en $x=l$ ($N=1$):}
\[
  \omega_1 = 2\omega_0\sen\!\left(\frac{\pi}{4}\right) = \omega_0\sqrt{2},
  \qquad \omega_0 = \sqrt{\frac{T}{ml}}
\]

\textit{(b) Dos masas iguales $m$ en $x=l$ y $x=2l$ ($N=2$, longitud total $3l$):}
\[
  \omega_1 = 2\omega_0\sen\!\left(\frac{\pi}{6}\right) = \omega_0
  \quad\text{(modo inferior)}
\]
\[
  \omega_2 = 2\omega_0\sen\!\left(\frac{2\pi}{6}\right) = \omega_0\sqrt{3}
  \quad\text{(modo superior)}
\]
En el modo inferior ambas masas se desplazan en el mismo sentido; en el modo superior en
sentidos opuestos (nodo en el centro).

\textit{(c) La frecuencia más alta es $\omega_2 = \omega_0\sqrt{3}$.}
\end{ejemplo}

\begin{ejemplo}[Osciladores acoplados forzados: determinación de \texorpdfstring{$X_1(t)$}{X1} y \texorpdfstring{$X_2(t)$}{X2}]
Considérese el sistema de dos masas $m$ unidas entre sí y a una pared por resortes de
constante $k$. La pared izquierda se mueve según $\eta(t) = \eta_0\cos\omega t$.

Las ecuaciones de movimiento son:
\begin{align*}
  m\ddot{x}_1 &= -k(x_1 - \eta) - k(x_1-x_2) = k\eta - 2kx_1 + kx_2 \\
  m\ddot{x}_2 &= -k(x_2 - x_1) = k(x_1-x_2)
\end{align*}

En coordenadas normales $q_1 = x_1+x_2$, $q_2 = x_1-x_2$ y con $\omega_0^2 = k/m$:
\[
  \ddot{q}_1 + \omega_0^2 q_1 = \omega_0^2\eta_0\cos\omega t,\qquad
  \ddot{q}_2 + 3\omega_0^2 q_2 = \omega_0^2\eta_0\cos\omega t
\]

Las soluciones estacionarias:
\[
  q_1 = \frac{\omega_0^2\eta_0}{\omega_0^2-\omega^2}\cos\omega t,\qquad
  q_2 = \frac{\omega_0^2\eta_0}{3\omega_0^2-\omega^2}\cos\omega t
\]

Recuperando las coordenadas originales: $x_1 = \frac{1}{2}(q_1+q_2)$,
$x_2 = \frac{1}{2}(q_1-q_2)$, con resonancias en $\omega = \omega_0$ y
$\omega = \omega_0\sqrt{3}$.
\end{ejemplo}

\end{document}
